% Preámbulo común a todas las figuras de los exámenes resueltos.
% Cada figura en _fuentes/figuras/<examen>/<nombre>.tex contiene solo el
% cuerpo (uno o más entornos tikzpicture / circuitikz); build.py lo envuelve
% con este preámbulo y lo compila con LuaLaTeX.
\documentclass[border=6pt,varwidth=40cm]{standalone}
\usepackage{fontspec}
\usepackage{unicode-math}
\setmainfont{Libertinus Sans}
\setmathfont{Libertinus Math}
\usepackage{xcolor}
\usepackage[siunitx]{circuitikz}
\usepackage{karnaugh-map}
\usepackage{tikz-timing}
\usepackage{array}
\usetikzlibrary{positioning,calc,arrows.meta,fit,backgrounds,shapes.geometric,
  decorations.pathreplacing,matrix}

% Paleta (la misma en todas las figuras)
\definecolor{tinta}{HTML}{1F2933}
\definecolor{acento}{HTML}{1E5AA8}
\definecolor{acentosuave}{HTML}{DCE8F7}
\definecolor{critico}{HTML}{C2410C}
\definecolor{criticosuave}{HTML}{FDE7DA}
\definecolor{ok}{HTML}{15803D}
\definecolor{oksuave}{HTML}{DCF2E3}
\definecolor{gris}{HTML}{9AA5B1}
\definecolor{grissuave}{HTML}{F1F3F5}

\ctikzset{logic ports=ieee, logic ports/scale=0.8, flipflops/scale=0.75}
\tikzset{
  >={Stealth[length=5pt]},
  every picture/.style={line width=0.7pt, draw=tinta, text=tinta},
  bloque/.style={draw, rounded corners=1.5pt, fill=white, minimum height=9mm,
    minimum width=16mm, align=center, font=\small},
  usado/.style={fill=acentosuave, draw=acento, line width=1pt},
  libre/.style={fill=grissuave, draw=gris, text=gris},
  rcritico/.style={draw=critico, line width=1.8pt},
  bus/.style={line width=1.6pt},
  nota/.style={font=\footnotesize\itshape, text=acento, align=left},
  etq/.style={font=\footnotesize},
}

% Multiplexor 2 a 1 dibujado como trapecio. #1 = nombre, #2 = posición,
% #3 = estilo extra (usado / libre).
\newcommand{\muxdos}[3]{%
  \node[draw, trapezium, shape border rotate=270, trapezium stretches body,
    minimum width=10mm, minimum height=5mm, trapezium angle=70,
    fill=white, #3, font=\scriptsize] (#1) at (#2) {};
  \node[font=\tiny, anchor=west] at ($(#1.north west)!0.28!(#1.south west)+(0.5mm,0)$) {1};
  \node[font=\tiny, anchor=west] at ($(#1.north west)!0.72!(#1.south west)+(0.5mm,0)$) {0};
}

% ---------------------------------------------------------------------------
% Slice simplificado de FPGA (Xilinx) tal como se ve en la presentación
% "Síntesis lógica": 4 LUT-6, dos MUX 2:1 (F7) y un MUX 2:1 (F8).
%
% \slicecfg{L4}{estilo L4}{estilo F7a}{estilo F7b}{estilo F8}
% \slice{x,y}{prefijo}{L1}{L2}{L3}{estilo L1}{estilo L2}{estilo L3}
% Los textos L1..L4 se escriben dentro de cada LUT. Los estilos son "usado"
% o "libre" (TeX admite 9 argumentos, por eso la LUT 4 y los MUX van aparte). Crea nodos <prefijo>l1..l4, <prefijo>fa, <prefijo>fb, <prefijo>f8
% y las coordenadas <prefijo>o1..o4 (salida directa de cada LUT), <prefijo>oa,
% <prefijo>ob (salidas de los MUX F7), <prefijo>o8 (salida del MUX F8),
% <prefijo>sa, <prefijo>sb, <prefijo>s8 (entrada de selección de cada MUX).
% ---------------------------------------------------------------------------
\newcommand{\slice}[8]{%
  \begin{scope}[shift={(#1)}]
    \foreach \i/\yy in {1/0,2/-2.1,3/-4.4,4/-6.5}{
      \coordinate (#2c\i) at (0,\yy);
    }
    \node[bloque, minimum width=22mm, minimum height=15mm, #6] (#2l1) at (#2c1) {LUT-6\\[1pt]#3};
    \node[bloque, minimum width=22mm, minimum height=15mm, #7] (#2l2) at (#2c2) {LUT-6\\[1pt]#4};
    \node[bloque, minimum width=22mm, minimum height=15mm, #8] (#2l3) at (#2c3) {LUT-6\\[1pt]#5};
    \sliceresto{#2}
  \end{scope}
}
% Segunda mitad (TeX solo admite 9 argumentos): LUT 4 y multiplexores.
% Se configura antes con \slicecfg{L4}{estilo L4}{F7a}{F7b}{F8}.
\newcommand{\slicecfg}[5]{%
  \def\slLcuatro{#1}\def\slEcuatro{#2}\def\slFa{#3}\def\slFb{#4}\def\slFocho{#5}}
\slicecfg{}{}{}{}{}
\newcommand{\sliceresto}[1]{%
  \node[bloque, minimum width=22mm, minimum height=15mm, \slEcuatro] (#1l4) at (#1c4) {LUT-6\\[1pt]\slLcuatro};
  % entradas de cada LUT (6 líneas)
  \foreach \l in {l1,l2,l3,l4}{
    \foreach \k in {1,...,6}{
      \draw[gris, line width=0.5pt] ($(#1\l.north west)!{(\k-0.5)/6}!(#1\l.south west)$) -- ++(-0.45,0);
    }
  }
  % salidas directas
  \foreach \l/\o in {l1/o1,l2/o2,l3/o3,l4/o4}{
    \draw (#1\l.east) -- ++(5.0,0) coordinate (#1\o);
    \coordinate (#1\o b) at ($(#1\l.east)+(4.6,0)$);
  }
  % MUX F7 (pares) y F8
  \muxdos{#1fa}{$(#1l1.east)!0.5!(#1l2.east)+(1.5,0)$}{\slFa}
  \muxdos{#1fb}{$(#1l3.east)!0.5!(#1l4.east)+(1.5,0)$}{\slFb}
  \muxdos{#1f8}{$(#1fa.east)!0.5!(#1fb.east)+(1.4,0)$}{\slFocho}
  \draw ($(#1l1.east)+(0.6,0)$) |- ($(#1fa.north west)!0.28!(#1fa.south west)$);
  \draw ($(#1l2.east)+(0.6,0)$) |- ($(#1fa.north west)!0.72!(#1fa.south west)$);
  \draw ($(#1l3.east)+(0.6,0)$) |- ($(#1fb.north west)!0.28!(#1fb.south west)$);
  \draw ($(#1l4.east)+(0.6,0)$) |- ($(#1fb.north west)!0.72!(#1fb.south west)$);
  \fill ($(#1l1.east)+(0.6,0)$) circle (1.2pt) ($(#1l2.east)+(0.6,0)$) circle (1.2pt)
        ($(#1l3.east)+(0.6,0)$) circle (1.2pt) ($(#1l4.east)+(0.6,0)$) circle (1.2pt);
  \draw (#1fa.east) -- ++(0.25,0) |- ($(#1f8.north west)!0.28!(#1f8.south west)$);
  \draw (#1fb.east) -- ++(0.25,0) |- ($(#1f8.north west)!0.72!(#1f8.south west)$);
  \draw (#1f8.east) -- (#1f8.east -| #1o1) coordinate (#1o8);
  \draw (#1fa.east) -- (#1fa.east -| #1o1) coordinate (#1oa);
  \draw (#1fb.east) -- (#1fb.east -| #1o1) coordinate (#1ob);
  \draw (#1fa.north) -- ++(0,0.35) coordinate (#1sa);
  \draw (#1fb.north) -- ++(0,0.35) coordinate (#1sb);
  \draw (#1f8.north) -- ++(0,0.35) coordinate (#1s8);
  \begin{scope}[on background layer]
    \node[draw=gris, dashed, rounded corners=4pt, inner sep=7pt,
      fit=(#1l1)(#1l4)(#1f8)(#1o1b)(#1sa)] (#1marco) {};
  \end{scope}
}
